%%%%%%%%%%%%%%%%%%%%%%%%%%%%%%%%%%%%%%%%%%%%%%%%%%%%%%%%%%%%%%%%%
\chapter{Unit 4: Designing Interpretable Predictive Models}
\label{chap:unit4}
%%%%%%%%%%%%%%%%%%%%%%%%%%%%%%%%%%%%%%%%%%%%%%%%%%%%%%%%%%%%%%%%%

\textbf{Total Time: 5 hours} \\
\textbf{Instructional Time: 3.5 hours} \\
\textbf{Assessment Time: 1.5 hours}

This unit introduces the foundations of supervised machine learning models, with a focus on interpretability and communicating modeling decisions.

%================================================================
\section{Foundations of Predictive Modeling}
\label{sec:4.1}
%================================================================

\textbf{Time: 1.5 hours (Instructional: 70 minutes, Assessment: 20 minutes)}

\noindent This lesson introduces supervised learning through the workflow that students will use throughout the rest of the unit: selecting a prediction target, choosing features, fitting an interpretable baseline model, and evaluating it on held-out data. The emphasis is on understanding what the model is doing, not only on obtaining a good score.

\subsection{Learning Objectives}

\begin{enumerate}
    \item Distinguish supervised learning from other common data analysis tasks, and distinguish regression from classification.
    \item Build a linear regression model as an interpretable baseline and explain what its coefficients represent.
    \item Practice the workflow for performing a train-test split, training a model on training data, evaluating it on held-out test data, and understanding how cross-validation works in practice.
    \item Identify common risks in predictive modeling, including overfitting and data leakage.
\end{enumerate}

\subsection{Assessment Instrument (20 minutes)}

Students will build a baseline linear regression model using the Vital Statistics dataset. They will document the outcome variable, list the features used, report one test-set performance metric such as RMSE or MAE, and write a brief note identifying two ways data leakage or overfitting could occur in this workflow.

%================================================================
\section{Interpretable Feature Engineering and Pipelines}
\label{sec:4.2}
%================================================================

\textbf{Time: 1 hour (Instructional: 40 minutes, Assessment: 20 minutes)}

\noindent This lesson focuses on interpretable feature engineering: creating features that better reflect patterns in the data while preserving interpretability. Students will use plots and domain knowledge to motivate simple transformations, encodings, and interactions, and will organize these steps within a small \verb|sklearn| Pipeline.

\subsection{Learning Objectives}

\begin{enumerate}
    \item Use visualizations to spot patterns in the data that may motivate feature engineering, such as nonlinearity, thresholds, and group differences.
    \item Understand when to use simple transformations, interaction terms, and common encoding strategies such as one-hot encoding.
    \item Build a small preprocessing and modeling Pipeline in \verb|sklearn| and explain how the resulting features affect model interpretability.
\end{enumerate}

\subsection{Assessment Instrument (20 minutes)}

Students will continue working with the Vital Statistics dataset and will add two or three engineered features in a small \verb|sklearn| Pipeline. They will submit the pipeline structure and a short written justification for each added feature, including why it may improve model fit or interpretability.

%================================================================
\section{Feature Justification, Multicollinearity, and Model Diagnostics}
\label{sec:4.3}
%================================================================

\textbf{Time: 1.5 hours (Instructional: 60 minutes, Assessment: 30 minutes)}

\noindent Building upon Section~\ref{sec:4.2}, students will learn how to decide whether a feature should remain in a model. The emphasis is on practical diagnostics and justification using correlations, multicollinearity checks, variance inflation factors, coefficient stability, and residual plots.

\subsection{Learning Objectives}

\begin{enumerate}
    \item Use simple quantitative tools such as Pearson and Spearman correlation, multiple $R^2$, and variance inflation factor to identify useful, redundant, or unstable features and to recognize multicollinearity.
    \item Use residual plots and side-by-side model comparison to assess whether added complexity is justified.
    \item Explain why a feature was kept, transformed, or removed based on predictive performance, interpretability, and stability.
\end{enumerate}

\subsection{Assessment Instrument (30 minutes)}

Students will compare two candidate models, such as a baseline model and an expanded feature-engineered model. They will decide which model they would keep and justify the decision using at least three pieces of evidence drawn from performance metrics, multicollinearity checks, coefficient stability, or residual diagnostics.

%================================================================
\section{Model Reporting and Stakeholder Communication}
\label{sec:4.4}
%================================================================

\textbf{Time: 1 hour (Instructional: 30 minutes, Assessment: 30 minutes)}

\noindent This lesson focuses on communicating model performance and design decisions to external stakeholders. The goal is to turn the analytic work from Section~\ref{sec:4.3} into a clear written summary that reports model results responsibly and in a form that aligns with the broader reporting expectations introduced earlier in the curriculum.

\subsection{Learning Objectives}

\begin{enumerate}
    \item Understand how to report and compare interpretable predictive models using appropriate regression metrics such as MSE, RMSE, and MAE.
    \item Summarize model design choices, strengths, and limitations in language appropriate for a non-technical audience.
    \item Practice connecting model reporting to broader ideas of transparency and reproducibility, including the spirit of TRIPOD-style reporting.
\end{enumerate}

\subsection{Assessment Instrument (30 minutes)}

Students will write a short report for a stakeholder describing their final model. The report should identify the prediction target, summarize the features included, report one or two performance metrics, and explain at least one important limitation or caveat in plain language.

% Required packages (add to your preamble if not already there):
% \usepackage{longtable}
% \usepackage{booktabs}
% \usepackage{ragged2e}
% \usepackage{array}
% \usepackage{pdflscape}

\begin{landscape}

    \section{Evaluation Rubric}

    {\footnotesize

    \begin{longtable}{%
        >{\RaggedRight\arraybackslash}p{2.5cm}   % Component
        >{\RaggedRight\arraybackslash}p{4.5cm}   % Excellent
        >{\RaggedRight\arraybackslash}p{4.2cm}   % Good
        >{\RaggedRight\arraybackslash}p{4.2cm}   % Satisfactory
        >{\RaggedRight\arraybackslash}p{4.5cm}}  % Needs Improvement

    \toprule
    \textbf{Component} &
    \textbf{Excellent (90--100\%)} &
    \textbf{Good (80--89\%)} &
    \textbf{Satisfactory (70--79\%)} &
    \textbf{Needs Improvement ($<$70\%)} \\
    \midrule
    \endfirsthead

    \toprule
    \textbf{Component} &
    \textbf{Excellent (90--100\%)} &
    \textbf{Good (80--89\%)} &
    \textbf{Satisfactory (70--79\%)} &
    \textbf{Needs Improvement ($<$70\%)} \\
    \midrule
    \endhead

    \midrule
    \multicolumn{5}{r}{\textit{Continued on next page}} \\
    \endfoot

    \bottomrule
    \endlastfoot

    % Row 1
    \textbf{1.\ Foundations of Predictive Modeling}
    \newline\textbf{(25\%)}
    &
    Clearly distinguishes supervised vs.\ unsupervised, regression
    vs.\ classification; correctly performs a train-test split, trains an
    interpretable baseline model, evaluates it on held-out data, and explains
    how cross-validation works; demonstrates strong understanding of
    overfitting and data leakage.
    &
    Good workflow with minor errors in reasoning or terminology;
    overfitting/data leakage described but not deeply.
    &
    Workflow partially correct; inconsistent understanding of splits,
    evaluation metrics, or model fitting.
    &
    Incorrect or incomplete workflow; poor understanding of model evaluation,
    overfitting, or data leakage.
    \\
    \midrule

    % Row 2
    \textbf{2.\ Interpretable Feature Engineering \& Pipelines}
    \newline\textbf{(20\%)}
    &
    Insightfully uses visualizations to identify meaningful patterns;
    appropriately applies simple transformations, encoding strategies, and
    interaction terms; builds a clear \verb|sklearn| Pipeline; thoroughly documents
    the rationale for each feature added.
    &
    Good feature construction with minor justification gaps; Pipeline mostly
    correct; documentation adequate.
    &
    Basic feature creation with limited insight; Pipeline incomplete or
    justification superficial.
    &
    Poor or incorrect feature construction; missing Pipeline; unclear or
    unjustified decisions.
    \\
    \midrule

    % Row 3
    \textbf{3.\ Feature Justification, Multicollinearity \& Diagnostics}
    \newline\textbf{(25\%)}
    &
    Correctly applies feature justification tools such as correlation checks,
    multiple $R^2$, variance inflation factor, and residual diagnostics;
    correctly identifies multicollinearity concerns; compares models
    thoughtfully and gives a clear, evidence-based rationale for keeping the
    preferred model.
    &
    Most diagnostics applied correctly; model comparison reasonable but with
    minor gaps in interpretation or justification.
    &
    Partial or inconsistent use of diagnostics; model choice weakly
    justified.
    &
    Incorrect or missing diagnostic evidence; model choice unsupported or
    inconsistent with results.
    \\
    \midrule

    % Row 4
    \textbf{4.\ Model Reporting \& Stakeholder Communication}
    \newline\textbf{(20\%)}
    &
    Accurately computes and interprets regression metrics such as MSE, RMSE,
    and MAE; report is clear, well-structured, and communicates model design
    decisions, performance, and limitations appropriately for the intended
    audience.
    &
    Metrics computed correctly with minor errors; report mostly clear and
    appropriately framed.
    &
    Metrics included but inconsistently computed or interpreted; explanation
    limited or too technical.
    &
    Metrics missing or incorrect; report unclear, incomplete, or not aligned
    with results.
    \\
    \midrule

    % Row 5
    \textbf{5.\ Overall Professionalism, Clarity, and Documentation}
    \newline\textbf{(10\%)}
    &
    Writing polished and well-organized; notebooks/scripts clean,
    reproducible, and well-commented; outputs easy to interpret.
    &
    Mostly clear with minor structural or formatting issues; code readable.
    &
    Understandable but inconsistent structure or documentation.
    &
    Poorly organized, unclear writing, missing documentation, or sloppy code.
    \\

    \end{longtable}

    } % end \footnotesize

\end{landscape}
